% ════════════════════════════════════════════════════════════════════
% SECTION 8 — LIMITATIONS AND FUTURE WORK  (v3.1: post-revision state)
% ════════════════════════════════════════════════════════════════════
\section{Limitations and Future Work}\label{sec:limitations}

\paragraph{Privacy scope.}
This work enforces the data-sovereignty constraint (raw data never
leaves the agency), which is the binding constraint of the
space-weather setting, but the stronger cryptographic notion of
privacy is not yet evaluated end-to-end. The communication transcript
of Eq.~\eqref{eq:privacy}---gradients or weight deltas---is known to
leak training information under gradient-inversion attacks
\cite{zhu2019}, particularly for over-parameterised models on
low-entropy data. The improved framework ships a
Bonawitz-style secure-aggregation module (verified by round-trip
tests) and a documented threat model, but the accuracy cost of
enabling secure aggregation and differential privacy on this benchmark
has not been measured; before any cross-agency deployment handling
genuinely sensitive observations, that budget must be re-measured.
Until then, ``data-locality-preserving'', not
``cryptographically private'', is the correct description of the
verified property.

\paragraph{Region-identity and validation inflation.}
The cleaned benchmark export does not carry active-region (HARP)
identifiers, so the stratified validation carve-out cannot enforce
region-level disjointness: the same active region may contribute
windows to both training and validation. This is why validation
metrics sit near ceiling (AUC $\approx 0.99$) while test metrics are
materially lower---the validation set is optimistically biased by
region overlap, though the dataset-level train/test boundary is
time-based and respected. The leakage audit reports this explicitly
rather than silently passing. The definitive fix is to re-partition by
active region using raw SWAN-SF metadata; the partitioning module
already accepts region identifiers when they are available.

\paragraph{Threshold transfer under prevalence shift.}
The frozen protocol selects $F_2$-optimal thresholds on validation
(48.9\% prevalence, mirroring the balanced training partitions), but
deployment prevalence is 1.9\%: the transferred thresholds put every
neural model in an almost-always-alert regime
(Section~\ref{sec:res-calib}). We deliberately did \emph{not} repair
this by touching test statistics---the v2.x evaluation effectively
did, and that is protocol leakage. The honest remedies are
prevalence-robust operating points (recall at fixed FPR, reported
here), prior correction using an \emph{estimated} deployment
prevalence from unlabelled data, or a natural-prevalence validation
carve-out. Until one of these is adopted, the threshold-conditional
test metrics (precision, recall, F1) should be read as documenting the
decision layer's failure mode, not as skill estimates.

\paragraph{Client-level evaluation semantics.}
The per-client comparison of Section~\ref{sec:res-clients} evaluates
each client's local model on a 20\% holdout of its own shard, but the
federated global models were trained on the \emph{full} shards,
including that holdout. The global models' per-client scores are
therefore within-distribution rather than clean generalisation
estimates. The comparison remains meaningful directionally (it
measures each client's experience of the global model on
client-distribution data), but absolute values are inflated for the
federated side; a region-disjoint client holdout is the rigorous
version.

\paragraph{Partition realism.}
Shard disjointness is now guaranteed by construction (index-level
allocation, audited at runtime), but the partition is still
\emph{statistical} label skew, not physical observatory coverage: no
mapping from active regions to the agencies whose instruments actually
observed them exists in the cleaned export. The partitioning module
supports a geographic mode that activates when region metadata is
available; until then, the six-client federation remains a simulation
of data locality, not a replication of it.

\paragraph{Algorithm coverage.}
The committed artefacts cover the MLP path; the LSTM path---which
consumes the raw $60\times24$ windows and matches the temporal
architectures that lead centralised SWAN-SF benchmarks
\cite{hassani2025}---is implemented but not yet evaluated federated,
and SCAFFOLD (implemented with control variates and crash recovery) is
disabled by default. Both are near-term experiments: the LSTM is the
most likely route to closing the gap to the tree-ensemble reference,
and SCAFFOLD is the natural third point on the drift-correction axis
between FedAvg and FedProx. Personalised FL (per-client heads over a
shared backbone) is a further direction: agencies may prefer a global
model for situational awareness plus a locally adapted model for their
own asset fleet.

\paragraph{Benchmark-construction effects.}
The cleaned variant's training partitions are pre-balanced (RUS,
Tomek links, TimeGAN augmentation), which (i) renders the SMOTE
balancing arm inert---reported as a null result---and (ii) creates
the $26\times$ prevalence shift that drives the calibration story. A
complementary evaluation on raw unbalanced SWAN-SF partitions would
separate benchmark-construction effects from federation effects, and
would allow the balancing ablation to actually engage.
